\documentclass[12pt,a4paper]{article}
\usepackage[utf8]{inputenc}
\usepackage[french]{babel}
\usepackage[T1]{fontenc}
\usepackage{amsmath,amssymb}
\usepackage{graphicx}
\usepackage{xcolor}
\usepackage{geometry}
\usepackage{tikz}
\usepackage{tcolorbox}
\usepackage{amsthm}
\geometry{margin=2.5cm}

% Définition des couleurs
\definecolor{titrebleu}{RGB}{0,102,204}
\definecolor{defvert}{RGB}{0,153,76}
\definecolor{exempleorange}{RGB}{255,102,0}
\definecolor{fondgris}{RGB}{245,245,245}
\definecolor{fondvert}{RGB}{230,255,230}
\definecolor{fondbleu}{RGB}{230,240,255}

% Boîtes colorées
\newtcolorbox{correction}{
    colback=fondvert,
    colframe=defvert,
    fonttitle=\bfseries,
    title=✓ Correction,
    arc=3mm
}

\newtcolorbox{demonstration}{
    colback=fondbleu,
    colframe=titrebleu,
    fonttitle=\bfseries,
    title=�� Démonstration,
    arc=3mm
}

\newtcolorbox{remarque}{
    colback=white,
    colframe=exempleorange,
    fonttitle=\bfseries,
    title=�� Remarque,
    arc=3mm
}

\title{\textbf{\Huge\color{titrebleu}Correction TD}\\[0.5em]
\large\color{defvert}Induction Mathématique}
\author{\large Olivier Raynaud}
\date{}

\begin{document}

\maketitle
\tableofcontents
\newpage

\section{Questions Préliminaires}

\subsection{Question 1 : Inégalité exponentielle}

\textbf{Énoncé :} Montrer par induction que si $x > 4$ alors $2^x > x^2$.

\begin{demonstration}
\textbf{Démonstration par récurrence sur $x \in \mathbb{N}$, $x > 4$ :}

\textbf{Cas de base :} Pour $x = 5$
\begin{itemize}
    \item $2^5 = 32$
    \item $5^2 = 25$
    \item On a bien $32 > 25$ ✓
\end{itemize}

\textbf{Hypothèse de récurrence :} Supposons que pour un certain $n \geq 5$, on ait $2^n > n^2$.

\textbf{Hérédité :} Montrons que $2^{n+1} > (n+1)^2$.

\begin{align*}
2^{n+1} &= 2 \cdot 2^n \\
&> 2 \cdot n^2 \quad \text{(par hypothèse de récurrence)} \\
&= n^2 + n^2
\end{align*}

Il suffit de montrer que $n^2 \geq (n+1)^2 - n^2 = 2n + 1$ pour $n \geq 5$.

\begin{align*}
n^2 &\geq 2n + 1 \\
n^2 - 2n - 1 &\geq 0
\end{align*}

Pour $n = 5$ : $25 - 10 - 1 = 14 > 0$ ✓

La fonction $f(n) = n^2 - 2n - 1$ est croissante pour $n > 1$, donc vraie pour tout $n \geq 5$.

Donc $2^{n+1} > n^2 + (2n+1) = (n+1)^2$ ✓

\textbf{Conclusion :} Par récurrence, $2^x > x^2$ pour tout $x \in \mathbb{N}$, $x > 4$.
\end{demonstration}

\subsection{Question 2 : Somme des carrés}

\textbf{Énoncé :} Démontrer par induction que $\displaystyle\sum_{i=0}^{n} i^2 = \frac{1}{6}n(n+1)(2n+1)$.

\begin{demonstration}
\textbf{Démonstration par récurrence :}

\textbf{Cas de base :} Pour $n = 0$
\begin{itemize}
    \item Membre de gauche : $0^2 = 0$
    \item Membre de droite : $\frac{1}{6} \cdot 0 \cdot 1 \cdot 1 = 0$
    \item Égalité vérifiée ✓
\end{itemize}

\textbf{Hypothèse de récurrence :} Supposons que pour un certain $n \geq 0$ :
$$\sum_{i=0}^{n} i^2 = \frac{1}{6}n(n+1)(2n+1)$$

\textbf{Hérédité :} Montrons que la formule est vraie pour $n+1$.

\begin{align*}
\sum_{i=0}^{n+1} i^2 &= \sum_{i=0}^{n} i^2 + (n+1)^2 \\
&= \frac{1}{6}n(n+1)(2n+1) + (n+1)^2 \quad \text{(par HR)} \\
&= (n+1)\left[\frac{n(2n+1)}{6} + (n+1)\right] \\
&= (n+1)\left[\frac{n(2n+1) + 6(n+1)}{6}\right] \\
&= (n+1)\left[\frac{2n^2 + n + 6n + 6}{6}\right] \\
&= (n+1)\left[\frac{2n^2 + 7n + 6}{6}\right] \\
&= (n+1)\left[\frac{(n+2)(2n+3)}{6}\right] \\
&= \frac{(n+1)(n+2)(2n+3)}{6} \\
&= \frac{(n+1)((n+1)+1)(2(n+1)+1)}{6} \quad ✓
\end{align*}

\textbf{Conclusion :} La formule est vraie pour tout $n \in \mathbb{N}$.
\end{demonstration}

\newpage
\section{Exercice 1 : Classe des arbres}

\subsection{Question 1 : Définition inductive des arbres}

\begin{correction}
\textbf{Définition inductive de la classe des arbres :}

Un arbre est défini inductivement par :
\begin{enumerate}
    \item \textbf{Cas de base :} Un sommet isolé est un arbre (avec 0 arête)
    \item \textbf{Cas inductif :} Si $T_1, T_2, \ldots, T_k$ sont des arbres, alors on peut former un nouvel arbre en créant un nouveau sommet racine et en le reliant aux racines de $T_1, \ldots, T_k$
\end{enumerate}
\end{correction}

\begin{demonstration}
\textbf{Propriété :} Tout arbre a un nœud de plus que son nombre d'arêtes.

\textbf{Preuve par induction structurelle :}

\textbf{Cas de base :} Un sommet isolé a 1 nœud et 0 arête. \\
On a bien $1 = 0 + 1$ ✓

\textbf{Hypothèse de récurrence :} Supposons que pour des arbres $T_1, \ldots, T_k$ avec respectivement $n_i$ nœuds et $a_i$ arêtes, on ait $n_i = a_i + 1$ pour tout $i$.

\textbf{Hérédité :} Construisons un nouvel arbre $T$ en ajoutant une racine reliée aux $k$ arbres.

\begin{itemize}
    \item Nombre de nœuds de $T$ : $n = 1 + \sum_{i=1}^{k} n_i$
    \item Nombre d'arêtes de $T$ : $a = k + \sum_{i=1}^{k} a_i$ (on ajoute $k$ arêtes)
\end{itemize}

Vérifions :
\begin{align*}
n &= 1 + \sum_{i=1}^{k} n_i \\
&= 1 + \sum_{i=1}^{k} (a_i + 1) \quad \text{(par HR)} \\
&= 1 + \sum_{i=1}^{k} a_i + k \\
&= \left(k + \sum_{i=1}^{k} a_i\right) + 1 \\
&= a + 1 \quad ✓
\end{align*}

\textbf{Conclusion :} Tout arbre a exactement un nœud de plus que son nombre d'arêtes.
\end{demonstration}

\subsection{Question 2 : Arbres binaires parfaits}

\begin{correction}
\textbf{Définition inductive des arbres binaires parfaits :}

\begin{enumerate}
    \item \textbf{Cas de base :} Une feuille isolée est un arbre binaire parfait de hauteur 0
    \item \textbf{Cas inductif :} Si $G$ et $D$ sont deux arbres binaires parfaits de même hauteur $h$, alors l'arbre formé d'une racine avec $G$ comme sous-arbre gauche et $D$ comme sous-arbre droit est un arbre binaire parfait de hauteur $h+1$
\end{enumerate}
\end{correction}

\subsection{Question 3 : Taille d'un arbre binaire parfait}

\begin{demonstration}
\textbf{Propriété :} Pour tout arbre binaire parfait $t$ de hauteur $h$, on a $|t| = 2^{h+1} - 1$.

\textbf{Preuve par induction structurelle :}

\textbf{Cas de base :} Une feuille a hauteur $h = 0$ et $|t| = 1$.
$$2^{0+1} - 1 = 2 - 1 = 1 \quad ✓$$

\textbf{Hypothèse de récurrence :} Supposons que pour deux arbres binaires parfaits $G$ et $D$ de hauteur $h$, on ait $|G| = |D| = 2^{h+1} - 1$.

\textbf{Hérédité :} Soit $T$ l'arbre binaire parfait de hauteur $h+1$ formé d'une racine avec sous-arbres $G$ et $D$.

\begin{align*}
|T| &= 1 + |G| + |D| \\
&= 1 + (2^{h+1} - 1) + (2^{h+1} - 1) \quad \text{(par HR)} \\
&= 1 + 2 \cdot 2^{h+1} - 2 \\
&= 2^{h+2} - 1 \\
&= 2^{(h+1)+1} - 1 \quad ✓
\end{align*}

\textbf{Conclusion :} La formule est vraie pour tout arbre binaire parfait.
\end{demonstration}

\subsection{Question 4 : Arbre syntaxique et grammaire de Chomsky}

\begin{demonstration}
\textbf{Propriété :} Si $T$ est un arbre syntaxique de hauteur $h$ produisant un mot $m$ à partir d'une grammaire en forme normale de Chomsky, alors $|m| \leq 2^{h-1}$.

\textbf{Rappel :} Une grammaire en forme normale de Chomsky n'a que des productions de la forme :
\begin{itemize}
    \item $A \to BC$ (deux non-terminaux)
    \item $A \to a$ (un terminal)
\end{itemize}

\textbf{Preuve par induction sur la hauteur :}

\textbf{Cas de base :} $h = 1$ (arbre réduit à une feuille)
\begin{itemize}
    \item L'arbre produit un seul symbole terminal : $|m| = 1$
    \item $2^{1-1} = 2^0 = 1$
    \item On a bien $1 \leq 1$ ✓
\end{itemize}

\textbf{Hypothèse de récurrence :} Pour tout arbre de hauteur $\leq h$, le mot produit a une longueur $\leq 2^{h-1}$.

\textbf{Hérédité :} Soit $T$ un arbre de hauteur $h+1$.

La racine utilise une production $S \to AB$. Les sous-arbres gauche et droit ont des hauteurs $\leq h$.

\begin{itemize}
    \item Soit $m_1$ le mot produit par le sous-arbre gauche : $|m_1| \leq 2^{h-1}$ (par HR)
    \item Soit $m_2$ le mot produit par le sous-arbre droit : $|m_2| \leq 2^{h-1}$ (par HR)
    \item Le mot produit est $m = m_1 m_2$
\end{itemize}

Donc :
\begin{align*}
|m| &= |m_1| + |m_2| \\
&\leq 2^{h-1} + 2^{h-1} \\
&= 2 \cdot 2^{h-1} \\
&= 2^h \\
&= 2^{(h+1)-1} \quad ✓
\end{align*}

\textbf{Conclusion :} La propriété est vérifiée pour tout arbre syntaxique.
\end{demonstration}

\newpage
\section{Exercice 2 : Triangle de Pascal et Code de Gray}

\subsection{Question 1 : Existence du code de Gray}

\begin{demonstration}
\textbf{Propriété :} Le code de Gray $T(n)$ existe pour tout entier naturel $n$.

\textbf{Preuve par récurrence :}

\textbf{Cas de base :} $n = 1$

Le code de Gray pour 1 bit est : $T(1) = (0, 1)$
\begin{itemize}
    \item On a bien $2^1 = 2$ éléments
    \item $0$ et $1$ diffèrent par exactement 1 bit ✓
\end{itemize}

\textbf{Hypothèse de récurrence :} Supposons qu'il existe un code de Gray $T(n)$ pour $n$ bits, contenant $2^n$ nombres où deux nombres consécutifs diffèrent par exactement 1 bit.

\textbf{Hérédité :} Construisons $T(n+1)$.

\textbf{Construction :}
\begin{enumerate}
    \item Préfixer tous les éléments de $T(n)$ par $0$ : on obtient la première moitié
    \item Préfixer tous les éléments de $T(n)$ \textit{en ordre inverse} par $1$ : on obtient la deuxième moitié
\end{enumerate}

\textbf{Exemple pour $n=2$ :}
\begin{itemize}
    \item $T(2) = (00, 01, 11, 10)$
    \item Préfixe $0$ : $(000, 001, 011, 010)$
    \item Préfixe $1$ inversé : $(110, 111, 101, 100)$
    \item $T(3) = (000, 001, 011, 010, 110, 111, 101, 100)$
\end{itemize}

\textbf{Vérification :}
\begin{itemize}
    \item \textbf{Taille :} $|T(n+1)| = 2 \cdot |T(n)| = 2 \cdot 2^n = 2^{n+1}$ ✓
    \item \textbf{Dans chaque moitié :} Les éléments consécutifs diffèrent par 1 bit (par HR) ✓
    \item \textbf{À la jonction :} Le dernier de la première moitié est $0xxx...x$ et le premier de la deuxième moitié est $1xxx...x$ où les bits suivants sont identiques (car on inverse l'ordre). Ils diffèrent donc exactement par le premier bit ✓
    \item \textbf{Circularité :} Le dernier élément $(100...0)$ et le premier $(000...0)$ diffèrent par 1 bit ✓
\end{itemize}

\textbf{Conclusion :} Le code de Gray existe pour tout $n \in \mathbb{N}$.
\end{demonstration}

\newpage
\section{Exercice 3 : Pavages}

\subsection{Question 1 : Divisibilité par 3}

\begin{demonstration}
\textbf{Propriété :} $2^{2n} - 1$ est divisible par 3 pour tout $n \in \mathbb{N}$.

\textbf{Preuve par récurrence :}

\textbf{Cas de base :} $n = 0$
$$2^{2 \cdot 0} - 1 = 2^0 - 1 = 1 - 1 = 0 = 3 \cdot 0 \quad ✓$$

\textbf{Hypothèse de récurrence :} Supposons que $2^{2n} - 1 = 3k$ pour un certain $k \in \mathbb{Z}$.

\textbf{Hérédité :} Montrons que $2^{2(n+1)} - 1$ est divisible par 3.

\begin{align*}
2^{2(n+1)} - 1 &= 2^{2n+2} - 1 \\
&= 4 \cdot 2^{2n} - 1 \\
&= 4 \cdot 2^{2n} - 4 + 3 \\
&= 4(2^{2n} - 1) + 3 \\
&= 4 \cdot 3k + 3 \quad \text{(par HR)} \\
&= 3(4k + 1) \quad ✓
\end{align*}

Donc $2^{2(n+1)} - 1$ est divisible par 3.

\textbf{Conclusion :} $2^{2n} - 1 \equiv 0 \pmod{3}$ pour tout $n \in \mathbb{N}$.
\end{demonstration}

\subsection{Question 2 : Pavage par triominos}

\begin{remarque}
Un \textbf{triomino} est une pièce en forme de L couvrant exactement 3 cases.
\end{remarque}

\begin{demonstration}
\textbf{Propriété :} Tout carré $2^n \times 2^n$ avec une case retirée peut être pavé par des triominos, pour $n \geq 1$.

\textbf{Preuve par récurrence :}

\textbf{Cas de base :} $n = 1$

Un carré $2 \times 2$ a 4 cases. Si on en retire une, il reste 3 cases qui forment exactement un triomino ✓

\textbf{Hypothèse de récurrence :} Supposons que tout carré $2^n \times 2^n$ avec une case retirée peut être pavé par des triominos.

\textbf{Hérédité :} Considérons un carré $2^{n+1} \times 2^{n+1}$ avec une case retirée.

\textbf{Construction :}
\begin{enumerate}
    \item Diviser le carré en 4 quadrants de taille $2^n \times 2^n$
    \item Un de ces quadrants contient la case retirée
    \item Placer un triomino au centre du grand carré, couvrant une case de chacun des 3 autres quadrants
    \item Maintenant, chaque quadrant a exactement une case "retirée"
    \item Par hypothèse de récurrence, chaque quadrant peut être pavé ✓
\end{enumerate}

\textbf{Illustration pour $n=2$ :}
\begin{verbatim}
+-------+-------+
| Q1    | Q2    |
|    X  |X      |  X = case retirée ou couverte par triomino central
+-------+-------+
| Q3    | Q4    |
|X      |    ■  |  ■ = case initialement retirée
+-------+-------+
\end{verbatim}

\textbf{Conclusion :} Tout carré $2^n \times 2^n$ avec une case retirée peut être pavé par des triominos pour $n \geq 1$.
\end{demonstration}

\newpage
\section{Résumé des techniques d'induction}

\begin{tcolorbox}[colback=fondbleu,colframe=titrebleu,title=\textbf{Types d'induction utilisés}]

\begin{enumerate}
    \item \textbf{Induction simple :} Questions 1, 2, Exercice 3
    \begin{itemize}
        \item Cas de base : vérifier pour la plus petite valeur
        \item Hérédité : montrer que $P(n) \Rightarrow P(n+1)$
    \end{itemize}
    
    \item \textbf{Induction structurelle :} Exercice 1
    \begin{itemize}
        \item Cas de base : structures élémentaires
        \item Hérédité : montrer que si la propriété est vraie pour les sous-structures, elle l'est pour la structure construite
    \end{itemize}
    
    \item \textbf{Induction constructive :} Exercice 2, Exercice 3.2
    \begin{itemize}
        \item Construction explicite de l'objet pour $n+1$ à partir de l'objet pour $n$
        \item Vérification que la construction préserve les propriétés requises
    \end{itemize}
\end{enumerate}
\end{tcolorbox}

\end{document}